\documentclass[12pt, a4paper]{article}
\usepackage[utf8]{inputenc}
\usepackage[T1]{fontenc}
\usepackage{microtype}
\usepackage[round, authoryear]{natbib}
\usepackage{hyperref}
\usepackage[margin=2.5cm]{geometry}

\title{Developmental Motivation}
\date{}

\begin{document}

\maketitle

\section*{Developmental Motivation}

Both halves of the central hypothesis have well-studied analogues in human and prenatal
development, which we take as motivation for our hypotheses.

For the first half --- the emergence of internal belief representations under optimisation
pressure --- a large developmental literature documents that infants and young children
behave in ways that are broadly consistent with several candidate optimisation objectives.
Children preferentially seek evidence at an intermediate level of surprise \citep[the
Goldilocks effect;][]{kidd2012goldilocks}, select informative causal interventions over
uninformative ones \citep{lapidow2020informative}, track learning progress
\citep{poli2020infants}, and adapt their search strategies to the ecological structure of
the learning environment \citep{ruggeri2022ecological}. Predictive processing accounts
recast many of these observations within a unified framework in which action selection is
guided by expected free energy, expressed in terms of anticipated prediction error
\citep{clark2013whatever, friston2016active}, which infants have also been found to
minimise \citep{zhang2019prediction}.

What this body of work cannot resolve, however, is the nature of the internal
representations underlying these behaviours. Infants behave as if they are minimising
surprise or maximising information gain, but the optimisation objectives themselves --- and
the internal states that implement them --- are not directly accessible. The behavioural
signatures are consistent with many different internal architectures, and manipulating the
objective directly is not possible in a living organism.

For the second half of the hypothesis --- the self/world distinction --- there also exist
strong developmental corollaries. Within hours of birth, infants alter their sucking rhythm
to produce preferred sounds, demonstrating rapid learning of an action--outcome contingency
\citep{decasper2009lateralized}. Newborns will raise an arm into a light beam and attend to
the resulting hand shadow, an early case of acting to produce a sensory consequence
\citep{vandermeer1997limelight}. By three months, infants kick more vigorously when their
ankle is tethered to a mobile that responds to their movement than under yoked
non-contingent control \citep{rovee1969conjugate}, and generalise the same
contingency-learning principle across modalities and effectors \citep{lee2013contingent}.
The mechanism proposed across these studies is temporal contingency detection: the infant
learns to identify its own actions by their reliable, short-latency coupling to reafferent
feedback, distinguishing self-caused from world-caused signal by the temporal structure of
the action--outcome link. This is precisely the causal-sampling property we use as the
signature of self in our toy model and LLM analyses --- a self-emitted token is recognisable
as self-produced by being a sample from one's own prior output distribution, in the same
way that an infant's movement is recognisable as self-produced by its temporal coupling to
the motor intention that preceded it. However, like with the developmental evidence
presented in support of our first hypothesis, this evidence does also not directly address
the link between optimisation function to explore well and the emergence of causal internal
representations.

Our own recent work (see attached) offers what might be read as circumstantial evidence
that optimisation pressure on exploration may itself drive the emergence of causal internal
representations. In utero, as soon as foetal movement begins, we observe a rapid and
systematic increase in the structure of exploratory motor output: early vermicular movements
(around 7.5 weeks GA; \citealt{prechtl1985ultrasound}) give way to increasingly
fractionated, spatially selective limb and joint movements (such as general movements,
startles, and brief myoclonic limb twitches; \citealt{devries1982emergence,
piontelli2010conclusions, tiriac2014selfgenerated, lindemann2016spindle}) by 20 weeks
\citep{devries1982emergence}, consistent with a shift from coarse, metabolically cheap
sampling toward targeted, high-signal-to-noise exploration of the body's own sensorimotor
space. Crucially, this increase in structured exploration is contemporaneous with the
earliest evidence of a primitive self-model: by approximately 22 weeks GA foetal hand
trajectories toward the mouth follow more direct paths than trajectories toward the eyes,
which decelerate earlier \citep{zoia2007evidence}. This body-specific calibration implies a
mapping from motor commands to their expected sensory consequences on a self-structured
target space. The co-emergence of more structured exploration alongside a primitive causal
model of the body is exactly the pattern our hypothesis predicts, but it is a correlational
observation in a system whose internal computations are opaque. Whether the structure of
exploration itself drove the formation of the self-model, or whether both reflect a common
maturational programme, cannot be determined from developmental data alone.

\bigskip\noindent\rule{\textwidth}{0.4pt}\bigskip

Taken together, these observations motivate both halves of our hypothesis and suggest that
the co-development of structured exploration and a causal self-model may be a general
feature of learning systems operating under joint optimisation pressure on exploration and
belief update. At the same time, they underscore what developmental research alone cannot
provide: direct access to the optimisation objectives, the ability to vary the training
regime while holding everything else fixed, and causal interventions on the representations
themselves. We see our proposed work as both inspired by and contributing to this
developmental picture --- by investigating, in a system where the ground truth is computable
and the training objective is under experimental control, whether the pressures that
developmental research can only observe from the outside are sufficient to produce the
internal factorisations it predicts.

\bibliographystyle{plainnat}
\bibliography{devmotivation}

\end{document}
